% 第 11 章：验证与校核总览
% 本章文件由 \input 引入，不含导言区。
\section{验证与校核总览}\label{sec:validation}

本章集中汇总潮流计算全家族的验证资产：测试矩阵、算例目录、外部对照通道、
已知边界与运行方法。各求解器节末的“验证与校核”小节（如第~\ref{sec:newton}~章）
给出局部依据，本章给出全局地图与诚实口径。

\subsection{验证体系分层}
平台潮流验证分四层，层层判据不同、互为补充：
\begin{enumerate}
  \item \textbf{单元与回归层}（Catch2，100+ 目标、约 1250 用例）：自建小算例的
        解析断言、收敛性回归、选项开/关对照；
  \item \textbf{外部对照层}：MATPOWER 三档回归、OpenDSS dss\_capi 桥、
        GridLAB-D 快照、Julia/JPC 相域基线；
  \item \textbf{解析与差分层}：雅可比中心差分校验、独立致密残差复算、
        两节点解析解；
  \item \textbf{能力合同与端到端层}：HTTP 能力覆盖矩阵与 GUI 全链路冒烟。
\end{enumerate}
铁律沿用第~\ref{sec:overview}~章的诚实结果口径：近似、fallback、门控跳过、
模型覆盖不足必须显式声明，验证报告不得把 skip 计为 pass。

\subsection{潮流测试矩阵总表}
\srcpath{tests/} 下潮流相关测试文件（仅列实际存在者，判据以文件内
\texttt{SECTION}/\texttt{REQUIRE}/\texttt{CHECK} 实际断言为准）：

\begin{footnotesize}
\begin{longtable}{@{}L{3.9cm}L{3.2cm}L{3.5cm}L{4.4cm}@{}}
\toprule
测试文件 & 覆盖求解器 & 算例/对照 & 容差/判据\\
\midrule
\endhead
\srcpath{test\_matpower\_cases.cpp} & 统一 Newton（+FDPF 一致性、DC OPF） &
  MATPOWER 三档（见 \ref{subsec:val-cases} 节）&
  收敛 + $|V|\in[0.5,1.5]$、$|\delta|<60^\circ$、残差 $<10^{-4}$；Tier2 批收敛率 $\ge 85\%$\\
\srcpath{test\_powerflow\_performance.cpp} & Newton（\texttt{[slow]} 标签）&
  MATPOWER 全目录（排除 \texttt{contab}/\texttt{scenarios}）&
  独立致密残差复算 $<10^{-6}$；电压窗 $(0.5,1.5)$\\
\srcpath{test\_unified\_pf\_upgrade.cpp} & 耦合雅可比/增强方程/半光滑/全局化 &
  解析雅可比 vs 中心差分（$\varepsilon=10^{-6}$）&
  $\max$ 相对误差 $<10^{-3}$；VDC\_VAC 定点 $<10^{-6}$\\
\srcpath{test\_scaling\_regression.cpp} & Newton 缩放开/关 &
  case9/14/30/57/118 &
  均收敛；缩放/原始残差与条件数估计有限\\
\srcpath{test\_helm\_solver.cpp} & HELM（+FDPF 稀疏注入、组件负荷聚合）&
  两母线电压稳定算例 vs Newton &
  HELM 与 NR 的 $|V|,|\theta|$ 差 $<2\times10^{-5}$；FDPF 与 NR 差 $<2\sim3\times10^{-4}$\\
\srcpath{test\_voltage\_stability.cpp} & CPF 弧长延拓 &
  两母线 P--V 鼻点 + 组件负荷方向 &
  翻越鼻点、$\lambda_{\max}\in(0.25,0.5)$、下支 $\lambda$ 回落；方向量精确到 $10^{-12}$\\
\srcpath{test\_homotopy\_continuation.cpp} & 同伦延拓 &
  两母线重负荷 &
  $\lambda=1.0$ 抵达端点、接受步 $\ge3$、收敛\\
\srcpath{test\_newton\_krylov.cpp} & NK-GMRES / 求解器工厂 / SolverWorkspace &
  $3\times3$ 耦合块系统 + 两母线 &
  GMRES 步误差 $<10^{-10}$；六方法工厂均可解；工作区复位保布局\\
\srcpath{test\_hacdcpf.cpp} & 公共 API 全族（AC/DC/FDPF/自适应/孤岛/线性化）&
  自建 2 母线 + IEEE14 AC/DC case builder &
  收敛、结果尺寸、有限性\\
\srcpath{test\_advanced\_pf.cpp} & 孤岛检测/自适应/分布式松弛/PV--PQ 策略 &
  9 母线、case14/118/2383wp &
  岛数精确、参与因子归一、双路径收敛一致\\
\srcpath{test\_dc\_island\_reference.cpp} & DC 岛参考规划/自动提升 &
  自建混合微网 &
  $V_{dc}$ 钉住、提升记录并告警、真参考不提升\\
\srcpath{test\_converter\_coordination.cpp} & 换流器协调预校核（46 个用例）&
  规则目录 \srcpath{ACDC-*/DCDC-*/DCISLAND-*} &
  阻断/告警语义、\fld{converter\_model\_scope} 声明\\
\srcpath{test\_all\_components\_pf.cpp} & 29 类组件全链路投影+潮流 &
  all29 集成系统 &
  收敛 + 每类组件对结果有可测影响\\
\srcpath{test\_distribution\_pipeline.cpp} & 配网 BFS 管线 &
  CB/开关/VSC rich$\to$graph$\to$PF$\to$OPF &
  管线端到端收敛\\
\srcpath{test\_three\_phase\_pf.cpp} & 三相投影与平衡化 &
  自建 3 母线（平衡/不平衡/环网/多源） &
  $|V|\in[0.8,1.1]$、迭代 $<15$、投影功率守恒\\
\srcpath{test\_three\_phase\_nr.cpp} & 相域 NR（16 个用例）&
  平衡/不平衡/支线/环网/多源 &
  残差 $<$ tol、VUF $\approx 0$、逐相功率平衡\\
\srcpath{test\_three\_phase\_nr\_load\_models.cpp} & 相域 NR 负荷模型 &
  wye/delta ZIP、中性点位移 &
  阻抗--电流--功率语义次序、不支持合同 fail-close\\
\srcpath{test\_three\_phase\_nr\_regulator.cpp} & 相域 NR 调压器闭环 &
  分相独立、远方 PT+LDC &
  档位单向提升、限值/振荡 fail-close\\
\srcpath{test\_three\_phase\_nr\_phase\_indexing.cpp} & 紧凑节点索引 &
  单相/两相支线 &
  无幻影相节点、phase\_mask 越界拒绝\\
\srcpath{test\_three\_phase\_nr\_transformer\_source.cpp} & 相域变压器源 &
  接地 wye 抽头、YNyn6 移相 &
  abc 域戳记精确\\
\srcpath{test\_three\_phase\_hybrid\_pf.cpp} & 相域 AC+DC 联立 &
  GFL Vdc--Q / GFM Vdc--内电势控制 &
  联立收敛（见第~\ref{sec:threephase}~章）\\
\srcpath{test\_opendss\_pf\_compare.cpp} & OpenDSS 桥对照（门控）&
  IEEE13、case33bw 全三相、调压器闭环 &
  分层 evidence/accepted 门槛，回归时退出码 2\\
\srcpath{test\_gridlabd\_compare.cpp} & GridLAB-D 快照对照 &
  两/三母线、元件阶梯、馈线阶梯 &
  运行时探测，缺失自动 skip\\
\srcpath{test\_phase\_domain\_ieee13.cpp} 等 plain 测试 & 相域求解器 vs JPC/OpenDSS 基线 &
  IEEE13/IEEE123 &
  \srcpath{HACDCPF\_HAVE\_OPENDSS} + 数据目录双重门控\\
\srcpath{test\_hacdcdss.cpp} & \srcpath{NetworkModel/PowerModels} 基元 &
  2 母线、IEEE14 AC/DC &
  $Y_{bus}$ 对称、注入自洽、收敛尺寸\\
\bottomrule
\end{longtable}
\end{footnotesize}

\subsection{算例目录}\label{subsec:val-cases}
\srcpath{external\_data/} 下潮流相关数据（均经 \texttt{ls} 核实）：
\begin{itemize}
  \item \textbf{\srcpath{matpower/}}：85 个 \texttt{.m} 文件，其中
        \texttt{contab\_*}（4 个）与 \texttt{scenarios\_*}（2 个）为扰动/场景
        辅助文件，潮流回归只取其余 79 个标准算例，分三档
        （\srcpath{tests/test\_matpower\_cases.cpp}）：Tier1 九算例
        （case5/9/14/30/57/118/300/39/24\_ieee\_rts）逐案断言收敛与电压合理性；
        Tier2 批量 54 算例（case533mt\_hi/lo 因解析器不支持算术表达式允许
        skip 不计失败）批收敛率 $\ge 85\%$；Tier3 大系统 14 算例至
        case\_SyntheticUSA（约 8.2 万节点，先关 PV--PQ 切换求解、失败再回退
        全开，:504--515）；
  \item \textbf{\srcpath{opendss\_ieee\_pes/}}：IEEE PES 测试馈线，
        \texttt{extracted/} 含 13/34/123 节点原始资料，
        \texttt{opendss\_reference/13\_node/official\_full/IEEE13Nodeckt.dss}
        是相域对照的官方基准（\srcpath{tests/CMakeLists.txt:250}）；
  \item \textbf{\srcpath{profiles/}}：\texttt{case33bw} 与
        \texttt{nansha\_network} 两套 8760 小时负荷/辐照/电价时序
        （时序生产模拟用）；
  \item \textbf{\srcpath{classical\_example/}}：4 个 JSON 算例（含
        \srcpath{hybrid\_ac\_dc\_mg\_case\_01.json}、
        \srcpath{nansha\_full\_network.json}）；
  \item \textbf{\srcpath{opf\_hand\_cases/}}：4 个手算 OPF 小算例（含 README
        给出解析期望；该目录当前不在仓库检出中）。
\end{itemize}

\subsection{外部对照通道细节}
\begin{itemize}
  \item \textbf{MATPOWER 回归}：Catch2 侧为上述三档测试；脚本侧
        \srcpath{tools/matpower\_benchmark.py}（217 行）驱动 MATLAB 批量
        \texttt{runpf} 与本仓库 \srcpath{tools/matpower\_pf\_compare.cpp}
        （输出收敛标志与逐母线电压 JSON，支持 \texttt{--no-pv-pq}/
        \texttt{--flows}/\texttt{--repeat N}），按匹配母线号比较
        $\max/\mathrm{mean}\,|\Delta V_m|$（pu）与 $|\Delta\theta|$（度），
        并对比两侧耗时；产物落在 \texttt{output/benchmarks/matpower\_pf/}。
        注意脚本内 MATLAB 路径与 \texttt{build/macos-release} 二进制路径为
        硬编码，跨机使用需改路径或走 Catch2 通道；
  \item \textbf{OpenDSS（dss\_capi 桥）}：三级门控——CMake 选项
        \srcpath{HACDCPF\_ENABLE\_OPENDSS}（默认 OFF）与
        \srcpath{HACDCPF\_ENABLE\_OPENDSS\_COMPARE} 先生成内部变量
        \srcpath{HACDCPF\_HAVE\_OPENDSS}（CMakeLists.txt:265--334），再派生
        \srcpath{HACDCPF\_HAVE\_OPENDSS\_COMPARE}（:solve\_once（lambda））；
        \srcpath{tests/test\_opendss\_pf\_compare.cpp}（3717 行）仅在后者为真时
        编译（\srcpath{tests/CMakeLists.txt:254--257}），覆盖单相/三相桥对照、
        case33bw 全三相 vs BFS 单相等值、调压器闭环与变压器分接头投影，
        接受层回归时对照门返回退出码 2；
  \item \textbf{GridLAB-D}：\srcpath{tests/test\_gridlabd\_compare.cpp} 目标
        始终构建，外部对照在运行时按 \texttt{HACDCPF\_GRIDLABD\_BIN}/
        \texttt{GRIDLABD\_BIN}/PATH/兄弟构建目录顺序探测可执行文件，缺失则
        干净 skip（\srcpath{tests/CMakeLists.txt:229--234}）；PV 母线电压调节
        不精确是已记录的 diagnostic 用例而非 pass；
  \item \textbf{解析解与雅可比中心差分}：
        \srcpath{src/power\_flow/jacobian\_check.cpp}（75 行）对任意
        残差/雅可比函数对做中心差分（默认 $h=10^{-6}$），按相对误差容差判定；
        \srcpath{src/power\_flow/residual\_evaluator.cpp}（157 行）从解向量
        以稀疏复数形式 $S=V\cdot\overline{Y_{bus}V}$ 独立重算注入残差
        （不再转致密矩阵，仅用于校核），被
        \srcpath{test\_powerflow\_performance.cpp} 以 $<10^{-6}$ 门槛复算
        每个收敛解（PV--PQ 切换或多岛算例因指定量不同显式跳过复算，:163--193）；
  \item \textbf{能力合同与 GUI E2E}：
        \srcpath{tools/validate\_case\_capabilities.py}（698 行）对 19 个能力域
        逐一发 HTTP 断言，输出 \texttt{output/capability\_coverage.md}
        覆盖矩阵（PASS/FAIL/SKIP 分列，skip 不计 pass）；
        \srcpath{tools/gui\_api\_e2e.py}（1050 行）启动
        \texttt{run\_gui\_server} 走 REST 全链路（装载算例$\to$ETAP 导出/再导入
        $\to$潮流$\to$短路），注册为 ctest 用例 \texttt{gui\_api\_e2e}
        （标签 \texttt{gui;e2e}、超时 120~s，由 \texttt{HACDCPF\_ENABLE\_ETAP}
        与 Python3 双重门控，\srcpath{tests/CMakeLists.txt:423--433}）。
\end{itemize}

\subsection{已知边界与覆盖缺口}
对照主文档“现状与诚实口径声明”逐条展开，每条附代码或测试证据：
\begin{enumerate}
  \item \textbf{HELM、同伦延拓、CPF、Newton--Krylov 已有专属回归（NEW 新增），
        但仍无生产 API 自动挂接}：实现分别在
        \srcpath{src/power\_flow/helm\_solver.cpp}、
        \srcpath{homotopy\_continuation.cpp}、\srcpath{voltage\_stability.cpp}、
        \srcpath{newton\_krylov.cpp}；NEW 起四者均有专属 Catch2 目标
        （\srcpath{tests/CMakeLists.txt:122--125}）：
        \srcpath{test\_helm\_solver.cpp}（HELM 与 Newton 在两母线算例上
        $|V|$/$|\theta|$ 一致至 $2\times10^{-5}$，另含 FDPF 稀疏注入与组件负荷
        聚合两条回归）、\srcpath{test\_homotopy\_continuation.cpp}
        （$\lambda$ 抵达 $1.0$ 端点）、\srcpath{test\_voltage\_stability.cpp}
        （弧长 CPF 翻越两母线 P--V 鼻点、比例方向计入组件负荷）、
        \srcpath{test\_newton\_krylov.cpp}（GMRES 块系统、求解器工厂、
        SolverWorkspace 复位）。覆盖深度仍属小算例 smoke 级，大系统回归
        仍缺；四者依旧由测试/GUI 直接驱动或作为选项层能力暴露。
        Newton--Krylov 回退默认关闭
        （\fld{enable\_newton\_krylov\_fallback\{false\}}，
        \srcpath{include/hacdcpf/power\_flow/globalization/robust\_nonlinear\_options.hpp:RobustNonlinearOptions}），
        仅在条件数代理越限时可选启用（src/power\_flow/newton\_solver.cpp:WorkspaceGuard）；
  \item \textbf{换流器物理不等式默认仅事后校核}：容量圆、交/直流电流限、调制比、
        占空比不嵌入 Newton 方程，求解后由
        \srcpath{src/api/hacdcpf.cpp:merge\_solver\_profiling} 发出
        \texttt{ACDC-PHYS-01..05}、:945 发出 \texttt{DCDC-PHYS-01} 告警；
        NEW 新增可选硬约束模式
        \fld{enforce\_converter\_physical\_limits\{false\}}
        （\srcpath{include/hacdcpf/power\_flow/power\_flow\_options.hpp:PowerFlowOptions}）：
        置 true 时带物理告警的 Newton 根被判不可行而非报收敛
        （\srcpath{src/api/hacdcpf.cpp:solve\_newton\_with\_lcc\_tap\_control}），并在结果 \fld{validity} 中
        如实标记各限值已强制（:solve\_newton\_with\_lcc\_tap\_control）；默认仍为仅告警模式。
        覆盖边界经 \fld{converter\_model\_scope} 声明（回归：
        \srcpath{test\_converter\_coordination.cpp} 的
        “Power-flow result declares its converter model scope”，:652）；
  \item \textbf{FDPF 注入已稀疏化}：\srcpath{src/power\_flow/fdpf\_solver.cpp}
        的注入计算改为稀疏复数形式 $S=V\cdot\overline{Y_{bus}V}$
        （\srcpath{compute\_power\_injections}，:19--35），不再转致密矩阵，
        原稠密 $O(n^2)$ 双重循环性能项随之消除；指定注入在
        \fld{has\_component\_loads} 时改用逐母线 ZIP 系数
        （\srcpath{compute\_specified\_injections}，:38--71）。回归：
        \srcpath{tests/test\_helm\_solver.cpp} 的 “FDPF sparse injection
        path remains numerically consistent” 与 “FDPF uses aggregated
        component loads”（与 NR 结果差 $<2\times10^{-4}$/$<3\times10^{-4}$）；
  \item \textbf{CPF 默认弧长延拓}：NEW 起 \srcpath{CpfSolver} 默认以弧长增广
        求解（\fld{enable\_arc\_length\{true\}}，
        \srcpath{include/hacdcpf/power\_flow/voltage\_stability.hpp:CpfOptions}），
        $\lambda$ 作为增广未知量，可翻越鼻点并保留下支
        \fld{lower\_branch\_steps\{2\}}（:CpfOptions）个点；置 false 时保留自然
        参数化旧模式（\srcpath{src/power\_flow/voltage\_stability.cpp:record（lambda）}）。
        \fld{CpfPoint} 新增 \fld{vdc}、\fld{va} 单位改为弧度（hpp:118--124）；
        鼻点摘要取迹线 $\lambda$ 最大点而非末点（cpp:495--500）。回归：
        \srcpath{tests/test\_voltage\_stability.cpp}（模型细节见
        第~\ref{sec:cpf}~章）；
  \item \textbf{\srcpath{workspace.hpp} 已落地为真实工作区}：
        \srcpath{include/hacdcpf/power\_flow/workspace.hpp} 新增
        \srcpath{prepare\_state}/\srcpath{prepare\_equations}/%
        \srcpath{reset\_iteration}（:35--42，实现
        \srcpath{src/power\_flow/workspace.cpp:prepare\_state}）；
        \srcpath{NewtonSolver} 复用每实例 \texttt{SolverWorkspace}，并以
        \fld{workspace\_in\_use\_} 递归守卫防止嵌套求解互踩
        （\srcpath{src/power\_flow/newton\_solver.cpp:build\_jacobian\_context}）。回归：
        \srcpath{tests/test\_newton\_krylov.cpp} 的 “SolverWorkspace resets
        values while retaining its layout”；
  \item \textbf{\srcpath{solver\_interface.hpp} 已接线}：NEW 新增
        \srcpath{src/power\_flow/solver\_interface.cpp}（137 行），实现
        \texttt{NewtonHybridSolver::solve} 等接口求解器与
        \srcpath{PowerFlowSolverFactory::create}
        （\srcpath{include/hacdcpf/power\_flow/solver\_factory.hpp:PowerFlowSolverFactory}）；
        工厂对 Newton/FDPF/DC/Adaptive/Continuation/NewtonKrylov 六方法
        均可产出可求解实例（回归 “Unified solver factory returns linked,
        callable solvers”，\srcpath{tests/test\_newton\_krylov.cpp:TEST\_CASE}）。
        生产 API 路径仍统一走 \srcpath{engine::NewtonSolver}
        （见第~\ref{sec:overview}~章求解器抽象层一节）。
\end{enumerate}

\subsection{如何运行验证}
configure/build/test 三段 preset（\srcpath{CMakePresets.json}，与根
\srcpath{AGENTS.md} 一致）：configure 侧 \texttt{macos-release}/
\texttt{linux-release}/\texttt{windows-msvc-release}/
\texttt{windows-vcpkg-release}/\texttt{full-dev}；test 侧同名 preset
仅前四个（\texttt{full-dev} 无 testPreset）。
\begin{quote}\footnotesize
\texttt{cmake --preset linux-release}\\
\texttt{cmake --build --preset linux-release}\\
\texttt{ctest --preset linux-release}
\end{quote}
常用定向运行（Catch2 标签与 ctest 正则）：
\begin{itemize}
  \item MATPOWER 回归：\texttt{ctest -R matpower\_cases}；慢速性能复算
        \texttt{test\_powerflow\_performance} 带 \texttt{[slow]} 标签，
        默认 Catch2 运行不含，需显式 \texttt{[slow]} 选择；
  \item 缩放/雅可比升级回归：\texttt{ctest -R}
        \srcpath{'scaling\_regression|unified\_pf\_upgrade'}；
  \item 非默认算法与工作区/工厂契约回归：\texttt{ctest -R}
        \srcpath{'helm\_solver|voltage\_stability|homotopy\_continuation|newton\_krylov'}；
  \item GUI 全链路：\texttt{ctest -R gui\_api\_e2e}（需
        \texttt{HACDCPF\_ENABLE\_ETAP=ON}，默认 ON，与 Python3）；
  \item OpenDSS 对照：configure 时加
        \texttt{-DHACDCPF\_ENABLE\_OPENDSS\_COMPARE=ON} 并提供 dss\_capi
        运行时；GridLAB-D 对照无需 configure 开关，运行时探测、缺失自动 skip；
  \item 能力合同矩阵：\texttt{python3}
        \srcpath{tools/validate\_case\_capabilities.py}
        \texttt{--server <run\_gui\_server 路径>}。
\end{itemize}
外部依赖（MATLAB、gridlabd、Julia、OpenDSS、chromium）缺失时对应测试自动
skip；判读报告时须区分“通过”与“跳过”，与主文档诚实口径一致。
